Když v hodnocení zohledníme dostupnost generativní AI, posun je z hodnocení produktu k hodnocení procesu, reflexe a doloženému přínosu studenta. Praktickým jádrem jsou tři prvky, které by každá rubrika měla obsahovat: pole pro deklaraci použití AI, metriky pro transparentnost a procesní důkazy, a disciplinární kritéria (správnost, argumentace, kódování apod.).

\medskip

Návrh základních AI‑aware položek do rubriky (doporučené škály)
\begin{itemize}
\item Transparentnost použití AI (0–5): deklarace nástrojů, promptů a stručné vysvětlení vlastního přínosu.
\item Originalita a vlastní přínos (0–10): úpravy, kritické zhodnocení a přidaná hodnota nad vygenerovaný obsah.
\item Doklad procesu (0–5): drafty, commit‑history, screencast nebo jiný procesní důkaz.
\end{itemize}

Škály vložte do rubriky a přiřaďte váhy podle typu úlohy.

\medskip

Praktické vzory rubrik pro tři typy úloh
\begin{enumerate}
\item Esej / písemná práce (příklad váh: Originalita 40\%, Transparentnost 20\%, Obsah 30\%, Styl 10\%)
  \begin{itemize}
  \item Originalita (0–10): co student upravil a proč; Transparentnost (0–5): nástroj/prompt/verze; Doklad procesu (0–5): drafty nebo historie; Obsah a argumentace (0–30).
  \end{itemize}

\item Programovací úloha / projekt (příklad: Správnost 50\%, Proces 30\%, Kód 20\%)
  \begin{itemize}
  \item Správnost (0–50): funkčnost podle testů; Doklad procesu (0–15): commit‑history, issue tracker, minimálně X commitů nebo screencast; Transparentnost (0–5); Kódová kvalita (0–30).
  \end{itemize}

\item Laboratorní zpráva / praktické cvičení (váhy přizpůsobit)
  \begin{itemize}
  \item Autentičnost dat (0–30): vlastní sběr vs. externí dataset; Procesní důkazy (0–30): protokoly, fotografie, surová data, skripty; Reflexe AI (0–10); Metodika a interpretace (0–30).
  \end{itemize}
\end{enumerate}

Implementační doporučení pro vyučující
\begin{itemize}
\item Do LMS vložte pole deklarace: checkbox o použití AI, krátké pole pro nástroj/verzi/prompt (150–300 znaků) a povinnou přílohu s procesním důkazem — zvyšuje transparentnost.
\item Pro kvantitativní testy využijte nástroje pro tvorbu a vyhodnocení kvízů, což usnadní konzistentní skórování \cite{kb_ai_101_tipu_pdf_page_36_1}.
\item U matematických cvičení nebo automatizovaných modulů vyžadujte nahrání postupu řešení (funkce „Pokrok v matematice“) \cite{kb_ai_101_tipu_pdf_page_8_1}.
\item Zvažte multimodální aktivity pro ověření porozumění a diskusi o AI (např. výukové světy v Minecraft Education) \cite{kb_ai_101_tipu_pdf_page_15_2}.
\end{itemize}

Krátké pokyny pro hodnotitele
\begin{itemize}
\item Před hodnocením zkontrolujte deklaraci a procesní přílohy; chybějící položky vyžádejte doplnit.
\item Při nejasnostech preferujte krátký dialog se studentem k verifikaci postupu a použijte rubriky pro formativní zpětnou vazbu — uveďte konkrétní kroky ke zlepšení.
\end{itemize}

Závěr: Přizpůsobené rubriky kombinují disciplinární kritéria s explicitními položkami pro AI (transparentnost, procesní důkazy, reflexe). Takový rámec snižuje motivaci k neohlášenému využití AI a posiluje kritické ověřování vygenerovaných výstupů.